% =============================================================================
%  §5  Ch.8 Monge–Mather 缩短原理
%  来源：Villani, Optimal Transport: Old and New, Chapter 8 (pp. 164–198)
%  主要定理：定理 8.1 / 推论 8.2（Mather 缩短引理），定理 8.5，定理 8.7
% =============================================================================
\TLsection{5 \ Monge--Mather 缩短原理}

% ── 动机（一）：最优传输轨迹不相交 ──────────────────────────────────────────
\begin{frame}{动机：最优传输轨迹不相交（Monge 观察）}
  \textbf{核心问题。} 如果两条最优运输路径相交，会发生什么？

  设 $c(x,y)=|x-y|^{2}$，$x_{1},x_{2}$ 是出发点，$y_{1},y_{2}$ 是目标点。

  \deflab{Monge 观察（Villani 第 8 章）。}
  \[
    |x_{1}-y_{1}|^{2}+|x_{2}-y_{2}|^{2} \;\le\; |x_{1}-y_{2}|^{2}+|x_{2}-y_{1}|^{2}.
  \]
  等价条件（展开后）：
  \[
    \langle x_{1}-x_{2},\,y_{1}-y_{2}\rangle \;\ge\; 0.
  \]

  \textbf{含义：} 相交赋值（$x_{1}\to y_{2}$，$x_{2}\to y_{1}$）的成本$\ge$不相交赋值（$x_{1}\to y_{1}$，$x_{2}\to y_{2}$）的成本。

  \textbf{推论：} 若存在更优的"不相交"配对，相交路径不可能是最优的——可以"解开"来降低总成本。

  \TLtakeaway{Monge 观察：对二次成本，最优赋值必为单调不相交赋值。交叉路径可通过"解开"严格降低总成本。}
\end{frame}

% ── 动机（二）：图示 ──────────────────────────────────────────────────────────
\begin{frame}{图示：相交 vs 不相交赋值（据 Villani 图 8.1 重绘）}
  \begin{center}
  \begin{tikzpicture}[scale=1.0, >={Stealth[length=5pt]}]
    \begin{scope}[shift={(-3.2,0)}]
      \node[font=\small,TLinkSoft] at (0,1.6) {\textbf{相交赋值}};
      \fill[TLprimary]  (-0.7, 0.9) circle (2.5pt) node[left,font=\scriptsize]{$x_1$};
      \fill[TLprimary]  ( 0.7, 0.9) circle (2.5pt) node[right,font=\scriptsize]{$x_2$};
      \fill[TLaccent]   (-0.7,-0.9) circle (2.5pt) node[left,font=\scriptsize]{$y_1$};
      \fill[TLaccent]   ( 0.7,-0.9) circle (2.5pt) node[right,font=\scriptsize]{$y_2$};
      \draw[->,TLprimary,thick] (-0.7, 0.9) -- ( 0.7,-0.9);
      \draw[->,TLaccent, thick] ( 0.7, 0.9) -- (-0.7,-0.9);
      \fill[TLneg] (0,0) circle (2.5pt);
      \node[TLneg,font=\scriptsize,right=1pt] at (0,0) {交叉};
    \end{scope}
    \draw[->,TLinkSoft,very thick] (-0.8,0) -- (0.8,0)
      node[midway,above,font=\scriptsize]{解开};
    \begin{scope}[shift={(3.2,0)}]
      \node[font=\small,TLinkSoft] at (0,1.6) {\textbf{不相交（最优）}};
      \fill[TLprimary]  (-0.7, 0.9) circle (2.5pt) node[left,font=\scriptsize]{$x_1$};
      \fill[TLprimary]  ( 0.7, 0.9) circle (2.5pt) node[right,font=\scriptsize]{$x_2$};
      \fill[TLaccent]   (-0.7,-0.9) circle (2.5pt) node[left,font=\scriptsize]{$y_1$};
      \fill[TLaccent]   ( 0.7,-0.9) circle (2.5pt) node[right,font=\scriptsize]{$y_2$};
      \draw[->,TLprimary,thick] (-0.7, 0.9) -- (-0.7,-0.9);
      \draw[->,TLaccent, thick] ( 0.7, 0.9) -- ( 0.7,-0.9);
      \node[TLpos,font=\scriptsize] at (0,-1.3) {成本更低};
    \end{scope}
  \end{tikzpicture}
  \end{center}
  {\scriptsize 据 Villani 重绘。}

  \textbf{数值验证（$x_1=0,\,x_2=1,\,y_1=0,\,y_2=1$）：}
  相交成本 $=|0-1|^2+|1-0|^2=2$；不相交成本 $=|0-0|^2+|1-1|^2=0$。节省了 $2$。

  \TLtakeaway{将相交赋值"解开"得到不相交赋值，总成本严格下降。这是第 8 章所有正则性结论的几何核心。}
\end{frame}

% ── 关键恒等式 (8.2) ──────────────────────────────────────────────────────────
\begin{frame}{关键恒等式：中间时刻控制全程（Villani 式 8.2）}
  \textbf{设置。} 令 $\gamma_{i}(t)=(1-t)x_{i}+ty_{i}$，$i=1,2$，$t\in[0,1]$。

  \deflab{恒等式 8.2（Villani）。}
  \begin{align*}
    &|\gamma_{1}(t)-\gamma_{2}(t)|^{2} \\
    &= (1-t)^{2}|x_{1}-x_{2}|^{2}+t^{2}|y_{1}-y_{2}|^{2} \\
    &\quad+t(1-t)\bigl[|x_{1}-y_{2}|^{2}+|x_{2}-y_{1}|^{2}-|x_{1}-y_{1}|^{2}-|x_{2}-y_{2}|^{2}\bigr].
  \end{align*}

  \textbf{若 $(x_1,y_1),(x_2,y_2)$ 满足最优性条件（括号内 $\ge 0$），则对任意 $t_0\in(0,1)$：}
  \[
    \sup_{0\le t\le 1}|\gamma_{1}(t)-\gamma_{2}(t)|
    \;\le\; \max\!\left(\tfrac{1}{t_{0}},\tfrac{1}{1-t_{0}}\right)
    \cdot|\gamma_{1}(t_{0})-\gamma_{2}(t_{0})|.
  \]

  \textbf{定性结论：} 若两最优轨迹在某中间时刻 $t_0$ 重合，则在整个 $[0,1]$ 上重合。

  \TLtakeaway{式 8.2 是 Monge 无交叉原理的精确定量化：中间时刻 $t_0$ 处的距离控制全程一致距离，乘以 $1/\min(t_0,1-t_0)$。这是第 8 章所有估计的代数基础。}
\end{frame}

% ── 定理 8.1（Mather 缩短引理）─────────────────────────────────────────────
\begin{frame}{定理 8.1：Mather 缩短引理（陈述）}
  \deflab{定理 8.1（Mather 缩短引理，Villani Thm 8.1）。}

  设 $M$ 为光滑 Riemann 流形，$d$ 为测地距离，$c(x,y)$ 由 Lagrange 量 $L(x,v,t)$ 定义。
  设 $x_{1},x_{2},y_{1},y_{2}\in M$ 满足 $c(x_{1},y_{1})+c(x_{2},y_{2})\le c(x_{1},y_{2})+c(x_{2},y_{1})$，
  $\gamma_{1},\gamma_{2}$ 为对应的作用量最小化曲线。假设：

  \begin{itemize}\small
    \item[(i)] 最小化曲线满足 Lipschitz 流（定义 7.6(d)）；
    \item[(ii)] $L$ 关于 $(x,v)$ 在 $\gamma_1,\gamma_2$ 的邻域中是 $C^{1,\alpha}$，$\alpha\in(0,1]$；
    \item[(iii)] $L$ 关于速度 $v$ 是 $(2+\kappa)$-凸的（严格凸的量化版本）。
  \end{itemize}

  则对任意 $t_{0}\in(0,1)$，存在常数 $C_{t_{0}}$ 和指数 $\beta=\beta(\alpha,\kappa)>0$ 使得
  \[
    \sup_{0\le t\le 1}d\!\bigl(\gamma_{1}(t),\gamma_{2}(t)\bigr)
    \;\le\; C_{t_{0}}\;d\!\bigl(\gamma_{1}(t_{0}),\gamma_{2}(t_{0})\bigr)^{\beta}.
  \]

  \TLtakeaway{定理 8.1（一般形式）：严格凸 Lagrange 量 $\Rightarrow$ 中间时刻距离 $\to$ 全程距离，指数 $\beta>0$。假设 (iii) 的 $(2+\kappa)$-凸性是量化的关键。}
\end{frame}

% ── 推论 8.2 ─────────────────────────────────────────────────────────────────
\begin{frame}{推论 8.2：$C^{2}$ Lagrange 量的缩短引理}
  \deflab{推论 8.2（Villani Cor 8.2）。}

  设 $L=L(x,v,t)$ 为 $C^{2}$ Lagrange 量，$\nabla_{v}^{2}L>0$（严格凸），$c$ 为对应成本。
  对任意紧集 $K\subset M$，存在 $C_{K}>0$：只要 $x_{1},y_{1},x_{2},y_{2}\in K$ 满足最优性条件，
  $\gamma_{1},\gamma_{2}$ 为对应的作用量最小化曲线，则对任意 $t_{0}\in(0,1)$，
  \[
    \sup_{0\le t\le 1}d\!\bigl(\gamma_{1}(t),\gamma_{2}(t)\bigr)
    \;\le\; \frac{C_{K}}{\min(t_{0},1-t_{0})}\;d\!\bigl(\gamma_{1}(t_{0}),\gamma_{2}(t_{0})\bigr).
  \]

  \textbf{最重要情形：} $c(x,y)=d(x,y)^{2}$（$L=|v|^{2}$，$\nabla_{v}^{2}L=2I>0$），推论 8.2 完整适用，指数 $\beta=1$。

  \textbf{例 8.4（Villani）。} $c=d^{1+\alpha}$（$0<\alpha<1$）：$L$ 非 $C^2$，但幂函数齐次性仍给出 $\beta=1$（见 Villani 附录）。

  % SOLUTION: 定理 8.1 完整证明含快捷路径构造和 Hölder 估计，见 Villani 第 172--180 页。留作 Part 6 备用幻灯片。

  \TLtakeaway{推论 8.2（最常用形式）：$C^2$、严格凸 Lagrange 量（含 $c=d^2$），$\sup d \le \frac{C_K}{\min(t_0,1-t_0)}\cdot d(t_0)$，指数 $\beta=1$。}
\end{frame}

% ── 定理 8.1 证明思路 ────────────────────────────────────────────────────────
\begin{frame}{定理 8.1 证明思路（概要）}
  \textbf{核心想法：} 若 $\gamma_{1},\gamma_{2}$ 在 $t_{0}$ 时刻相交于 $m_{0}$，且速度分别为 $v_{1},v_{2}$，就在短时间窗内构造"快捷路径"互换连接。

  \textbf{步骤一（局部化）。} 将流形问题化归为局部 Euclid 计算：紧性覆盖 $\gamma_{1}\cup\gamma_{2}$ 邻域，在小球中引入坐标。

  \deflab{补一个事实：$(2+\kappa)$-凸性。}
  这是 $L$ 关于速度的量化严格凸性：存在 $c_{\mathrm{gap}}>0$，使得
  \[
    L\!\left(x,\frac{v_{1}+v_{2}}{2},t\right)
    \le \frac{1}{2}L(x,v_{1},t)+\frac{1}{2}L(x,v_{2},t)
    -c_{\mathrm{gap}}\abs{v_{1}-v_{2}}^{2+\kappa}.
  \]
  这就是 Villani 使用的 midpoint-gap 形式。

  \TLtakeaway{定理 8.1 的局部证明只用两个工具：把曲线放进 Euclid 坐标，以及把速度严格凸性读成 midpoint-gap。}
\end{frame}

\begin{frame}{定理 8.1 证明思路（快捷路径主项）}
  \textbf{步骤二（快捷路径）。}
  在 $[t_{0}-\tau,t_{0}+\tau]$ 内，将两段曲线剪开并互换连接；两条新路径都近似以平均速度 $(v_{1}+v_{2})/2$ 穿过 $m_{0}$。主阶作用量改变为
  \[
    \Delta A \approx
    2\tau\!\left[2L\!\left(m_{0},\frac{v_{1}+v_{2}}{2},t_{0}\right)
    -L(m_{0},v_{1},t_{0})-L(m_{0},v_{2},t_{0})\right].
  \]

  \TLtakeaway{快捷路径为什么会降作用量：两条原速度 $v_{1},v_{2}$ 被两个平均速度替代，主项正是 midpoint-gap 的负数。}
\end{frame}

\begin{frame}{定理 8.1 证明思路（凸性矛盾）}
  \textbf{步骤三（严格下降）。}
  把上一帧的主项代入 midpoint-gap：若 $v_{1}\ne v_{2}$，则
  \[
    \Delta A
    \lesssim -4\tau c_{\mathrm{gap}}\abs{v_{1}-v_{2}}^{2+\kappa}<0.
  \]
  因而互换后的两条路径总作用量更小，即
  $c(x_{1},y_{2})+c(x_{2},y_{1})<c(x_{1},y_{1})+c(x_{2},y_{2})$，这与定理 8.1 的最优性假设矛盾。

  \textbf{步骤四（从相同速度到相同曲线）。}
  唯一剩下的可能是 $v_{1}=v_{2}$。此时两条曲线在 $t_{0}$ 处有同一位置 $m_{0}$ 和同一速度；最小化曲线满足 Lipschitz 流，故 Cauchy--Lipschitz 唯一性推出 $\gamma_{1}=\gamma_{2}$。

  \deflab{本节诚实省略的技术估计。}
  省略的是 Villani 第 172--180 页中，从快捷路径误差界 (8.21) 与凸性下界 (8.22) 推出速度 Hölder 控制 (8.8) 的计算；主线只需要知道该估计把"相交矛盾"量化成 (8.4)。
  % SOLUTION: 完整速度估计 (8.8) 的推导留作 Part 6 备用幻灯片。

  \TLtakeaway{定性结论来自严格凸性矛盾；定量结论来自 (8.21)--(8.22) 的误差比较，最终得到速度差被位置差的幂控制。}
\end{frame}

% ── 定理 8.5：陈述 ────────────────────────────────────────────────────────────
\begin{frame}{定理 8.5：中间时刻传输映射是 Lipschitz（陈述）}
  \deflab{定理 8.5（Villani Thm 8.5）。}

  设 $M$ 为 Riemann 流形，$c$ 满足推论 8.2 假设，$K\subset M$ 紧，$\Pi$ 为支撑在 $K$ 内的动力学最优传输。

  则 $\Pi$-几乎处处，对任意 $\gamma,\tilde{\gamma}\in\mathrm{supp}(\Pi)$，
  \[
    \sup_{0\le t\le 1}d\!\bigl(\gamma(t),\tilde{\gamma}(t)\bigr)
    \;\le\; \frac{C_{K}}{\min(t_{0},1-t_{0})}\,d\!\bigl(\gamma(t_{0}),\tilde{\gamma}(t_{0})\bigr).
  \]

  \textbf{关键推论（传输映射）。} 若 $(\mu_{t})$ 为紧支撑测度间的位移插值，$t_{0}\in(0,1)$，
  则映射 $T_{t_{0}\to t}\colon\gamma(t_{0})\mapsto\gamma(t)$ 在 $\mu_{t_{0}}$-a.s. 有定义，
  在定义域上是 Lipschitz 连续的（Lipschitz 常数 $C_{K}/\min(t_{0},1-t_{0})$），
  且是 $\mu_{t_{0}}$ 到 $\mu_{t}$ 的 Monge 问题的唯一解。

  \TLtakeaway{定理 8.5：中间时刻 $t_0\in(0,1)$ 的传输映射自动 Lipschitz，即使端点传输映射可能不正则——这是缩短引理的第一个重大应用。}
\end{frame}

% ── 定理 8.5：证明思路 ──────────────────────────────────────────────────────
\begin{frame}{定理 8.5 证明思路}
  \textbf{步骤一（最优性传递）。}
  $\pi=(e_{0},e_{1})_{\#}\Pi$ 为最优耦合，故 $c$-循环单调，因此 $\Pi\otimes\Pi$-a.s.：
  \[
    c(\gamma(0),\gamma(1))+c(\tilde{\gamma}(0),\tilde{\gamma}(1))
    \;\le\; c(\gamma(0),\tilde{\gamma}(1))+c(\tilde{\gamma}(0),\gamma(1)).
  \]

  \textbf{步骤二（推论 8.2）。} 上式即推论 8.2 前提，直接得到 Lipschitz 估计
  $\sup_{t}\,d(\gamma(t),\tilde{\gamma}(t))\le\frac{C_{K}}{\min(t_{0},1-t_{0})}\,d(\gamma(t_{0}),\tilde{\gamma}(t_{0}))$。

  \deflab{补一个事实：良定义就是单值。}
  若同一个中间点可由两条支撑曲线到达，即 $\gamma(t_{0})=\tilde{\gamma}(t_{0})$，则右端为 $0$。

  \textbf{步骤三（良定义与唯一性）。}
  缩短估计强制 $\sup_{t}d(\gamma(t),\tilde{\gamma}(t))=0$，所以 $\gamma=\tilde{\gamma}$。
  因而 $T_{t_{0}\to t}\colon\gamma(t_{0})\mapsto\gamma(t)$ 不依赖曲线选择，是单值良定义的；同一估计给出 Lipschitz 常数。最优唯一性由 §2 定理 5.30（Monge 可解性）给出。

  \TLtakeaway{$c$-循环单调 $\Rightarrow$ 推论 8.2 $\Rightarrow$ 缩短估计；中间点若相同，则整条曲线相同，所以 $T_{t_{0}\to t}$ 是单值良定义的。}
\end{frame}

\begin{frame}{例 8.6：中间映射为何显然 Lipschitz}
  \textbf{例 8.6（Villani）。} $c=|\cdot|^{2}$，$\mu_{0}=\delta_{0}$，$\mu=\law(X)$：
  随机测地线为 $\gamma(t)=t\gamma(1)$，所以
  \[
    \mu_{t}=\law(tX),
    \qquad
    T_{t_{0}\to t}=\frac{t}{t_{0}}\Id .
  \]
  该映射的 Lipschitz 常数是 $t/t_{0}$；这正是定理 8.5 在二次成本最简单情形下的显式模型。

  \TLtakeaway{例 8.6 把抽象结论化成公式：中间位置 $t_{0}X$ 唯一决定任意时刻 $tX$，映射就是线性缩放。}
\end{frame}

% ── 定理 8.7：陈述 ────────────────────────────────────────────────────────────
\begin{frame}{定理 8.7：位移插值保持绝对连续性（陈述）}
  \deflab{定理 8.7（绝对连续性，Villani Thm 8.7）。}

  设 $M$ 为 Riemann 流形，$L\in C^{2}$，$\nabla_{v}^{2}L>0$，$c$ 为对应成本。
  设 $\mu_{0},\mu_{1}\in\Pp(M)$，$C(\mu_{0},\mu_{1})<+\infty$，$(\mu_{t})_{0\le t\le 1}$ 为位移插值。

  \begin{center}
  \textbf{若 $\mu_{0}$ 或 $\mu_{1}$ 中至少一个绝对连续于 $M$ 上的体积元，}\\[4pt]
  \textbf{则 $\mu_{t}\in\Pac(M)$ 对所有 $t\in(0,1)$ 成立。}
  \end{center}
  {\small （此处 $\Pac$ 为绝对连续概率测度，定义见 §0。）}

  \textbf{直观。} 在中间时刻 $t_0$，传输映射 $T_{t_0\to 1}$ 是 Lipschitz 的（定理 8.5）。
  Lipschitz 映射的像保持零测集的零测性（Lusin N 条件），从而绝对连续性反向传递：
  $\mu_1\in\Pac$，$T_{t_0\to 1}$ Lipschitz，故 $\mu_{t_0}\in\Pac$。

  \TLtakeaway{定理 8.7：端点之一绝对连续 $\Rightarrow$ 中间时刻的插值测度自动绝对连续。内部正则性从最优传输的组合结构中"免费涌现"。}
\end{frame}

% ── 定理 8.7：证明思路 ──────────────────────────────────────────────────────
\begin{frame}{定理 8.7 证明（紧支撑情形）}
  \textbf{假设 $\mu_{1}\in\Pac$，$\mu_{0},\mu_{1}$ 紧支撑，证明 $\mu_{t_{0}}\in\Pac$，$t_{0}\in(0,1)$。}

  \textbf{步骤一。} 由定理 8.5（Lipschitz 性来自推论 8.2），存在 Lipschitz Monge 映射 $T\colon\mu_{t_{0}}\to\mu_{1}$。

  \deflab{补一个事实：Lusin (N) 性质 / 面积公式。}
  若 $T\colon\R^{n}\to\R^{n}$ 是 $L_{T}$-Lipschitz，则
  $\mathrm{vol}(T(N))\le L_{T}^{n}\mathrm{vol}(N)$；零测集仍送到零测集。流形上用局部坐标。

  \textbf{步骤二（零测集）。} 设 $N$ 为零体积集，并取 $L_{T}=C_{K}/\min(t_{0},1-t_{0})$。于是
  \[
    \mathrm{vol}(N)=0
    \;\Longrightarrow\;
    \mathrm{vol}(T(N))\le L_{T}^{n}\mathrm{vol}(N)=0
    \;\Longrightarrow\;
    \mu_{1}[T(N)]=0.
  \]
  最后一箭头用 $\mu_{1}\in\Pac$，见 §0。

  \TLtakeaway{Lusin (N) / 面积公式在这里只做一件事：$N$ 零体积 $\Rightarrow$ $T(N)$ 零体积 $\Rightarrow$ $\mu_{1}[T(N)]=0$。}
\end{frame}

\begin{frame}{定理 8.7 证明（紧支撑收尾）}
  \textbf{步骤三（绝对连续性）。} 由 $N\subset T^{-1}(T(N))$：
  \[
    \mu_{t_{0}}[N]\;\le\;\mu_{t_{0}}\!\left[T^{-1}(T(N))\right]
    =\pf{T}{\mu_{t_{0}}}[T(N)]=\mu_{1}[T(N)]=0,
  \]
  故 $\mu_{t_{0}}\in\Pac$。

  \textbf{可测性小修。} Villani 指出 $T(N)$ 未必是 Borel 集；严格写法是用一个零体积 Borel 集包含 $T(N)$，上式按该 Borel 包含集解释即可。

  \TLtakeaway{$T$ Lipschitz $+$ $\mu_{1}\in\Pac$ $\Rightarrow$ $T(N)$ 零测 $\Rightarrow$ $\mu_{t_{0}}[N]=0$：绝对连续性沿 Lipschitz 传输反向传递。}
\end{frame}

% ── 定理 8.7：证明续（非紧情形）──────────────────────────────────────────────
\begin{frame}{定理 8.7 证明续（非紧支撑情形）}
  \textbf{步骤四（非紧推广）。} 若 $\mu_{t_0}\notin\Pac$，取体积零但 $\mu_{t_0}$-正测度的集 $Z_{t_0}$；
  对 $\Pi$ 限制到紧子集 $K\subset\{\gamma\colon\gamma(t_0)\in Z_{t_0}\}$，化为上述紧支撑情形，得矛盾。

  \textbf{小结。} 三步骨架：
  \begin{enumerate}\small
    \item 定理 8.5 $\to$ $T$ Lipschitz（推论 8.2 缩短引理）。
    \item Lipschitz 面积公式 $\to$ $T(N)$ 零体积。
    \item $\mu_1\in\Pac$ $\to$ $\mu_1[T(N)]=0$ $\to$ $\mu_{t_0}[N]=0$。
  \end{enumerate}

  \TLtakeaway{绝对连续性沿 Lipschitz 传输的逆向传递：$\mu_1\in\Pac$（见 §0）$+$ $T$ Lipschitz $\Rightarrow$ $\mu_{t_0}\in\Pac$。非紧情形通过紧子集的穷举化归完成。}
\end{frame}

% ── 为什么定理 8.7 重要 ────────────────────────────────────────────────────
\begin{frame}{为什么定理 8.7 重要？内部正则性的来源}
  \textbf{看似悖论。} 端点 $\mu_0$ 可以非常奇异（原子测度、奇异测度），
  但只要 $\mu_1\in\Pac$，中间时刻 $\mu_t$ 就自动绝对连续。

  \begin{center}
  \begin{tikzpicture}[scale=0.9, >={Stealth[length=4pt]}]
    \node[draw=TLprimary,rounded corners=3pt,inner sep=5pt,font=\small,
          fill=TLprimaryTint] (L) at (0,0) {$\mu_0$（可奇异）};
    \node[draw=TLaccent,rounded corners=3pt,inner sep=5pt,font=\small,
          fill=TLaccentLight] (R) at (7.5,0) {$\mu_1\in\Pac$};
    \node[draw=TLpos,rounded corners=3pt,inner sep=5pt,font=\small,
          fill=TLpos!10] (M) at (3.75,0) {$\mu_t\in\Pac$，$\forall\,t\in(0,1)$};
    \draw[->,TLprimary,thick] (L) -- (M);
    \draw[->,TLaccent,thick] (R) -- (M);
  \end{tikzpicture}
  \end{center}

  \textbf{物理直觉。} 质量在中间时刻的位置由 Lipschitz 映射决定，不能聚焦到零测集。

  \textbf{整体逻辑链：}
  \[
    \underbrace{c\text{-循环单调}}_{\text{最优性}}
    \;\overset{\text{推论 8.2}}{\Rightarrow}\;
    \underbrace{T_{t_0\to t}\text{ Lipschitz}}_{\text{定理 8.5}}
    \;\overset{\text{Lusin N}}{\Rightarrow}\;
    \underbrace{\mu_t\in\Pac}_{\text{定理 8.7}}
  \]

  \textbf{前瞻。} 第 9–10 章以此内部绝对连续性为基础研究 Wasserstein 几何的曲率下界。

  \TLtakeaway{"内部正则性免费涌现"：最优传输的组合最优性（无交叉）通过量化（缩短引理）产生分析正则性（$\mu_t\in\Pac$）——这是最优传输几何的核心奇迹之一。}
\end{frame}

% ── 习题 5 · Exercises（一）────────────────────────────────────────────────
\begin{frame}{习题 5 · Exercises（一）}
  \diffI\ \textbf{习题 5.1.} 设 $x_1=0$，$x_2=2$，对 $c(x,y)=(x-y)^2$，
  比较赋值 $(x_1\to 0,\,x_2\to 2)$ 与 $(x_1\to 2,\,x_2\to 0)$ 的总成本，
  验证不相交赋值严格更优。

  \hint{不相交成本：$(0-0)^2+(2-2)^2=0$；相交成本：$(0-2)^2+(2-0)^2=8$。差值 $8=2\langle x_1-x_2,y_2-y_1\rangle$，与 $\langle x_1-x_2,y_1-y_2\rangle\ge 0$ 对应。}
  % SOLUTION: 不相交成本=0，相交成本=8，节省8。一般形式：展开差值得2(x1-x2)(y2-y1)>=0 当 x1<x2,y1<y2 时。

  \vspace{6pt}
  \diffII\ \textbf{习题 5.2.} 设 $(x_1,y_1),(x_2,y_2)\in\R\times\R$ 是某最优耦合
  $\pi$（成本 $c(x,y)=(x-y)^2$）支撑集中的两点。
  证明 $x_1<x_2\;\Rightarrow\;y_1\le y_2$（单调性）。

  \hint{利用 $\pi$ 的 $c$-循环单调性：$c(x_1,y_1)+c(x_2,y_2)\le c(x_1,y_2)+c(x_2,y_1)$。展开后等价于 $(x_1-x_2)(y_1-y_2)\ge 0$。$x_1<x_2$ 时 $x_1-x_2<0$，故 $y_1-y_2\le 0$。}
  % SOLUTION: c-循环单调给出(x1-y1)^2+(x2-y2)^2<=(x1-y2)^2+(x2-y1)^2。展开：(x1-x2)(y1-y2)>=0。x1<x2 时 y1<=y2。

  \TLtakeaway{习题 5.1 数值验证无交叉原理；习题 5.2 代数证明最优耦合支撑集的单调性——二者共同将 Monge 观察从直觉化为精确命题。}
\end{frame}

% ── 习题 5 · Exercises（二）────────────────────────────────────────────────
\begin{frame}{习题 5 · Exercises（二）}
  \diffIII\ \textbf{习题 5.3.} 设 $\mu_0\in\Pac(\R^n)$（绝对连续），$\mu_1$ 为任意概率测度，
  $c(x,y)=|x-y|^2$，$C(\mu_0,\mu_1)<+\infty$。
  利用定理 8.7 证明位移插值 $\mu_t$ 在 $t\in(0,1)$ 上绝对连续。

  \begin{itemize}
    \item[(a)] 识别定理 8.7 适用所需的假设，逐条验证。
    \item[(b)] $\mu_0\in\Pac$ 与 $\mu_1\in\Pac$ 在定理中的角色是否对称？
    \item[(c)] 若改为 $\mu_0$ 是原子测度（非绝对连续），$\mu_1\in\Pac$，结论是否仍成立？
  \end{itemize}

  \hint{(a) 检查：$L=|v|^2$ 是 $C^2$ 且 $\nabla_v^2L=2I>0$；成本有限；$\mu_0\in\Pac$ 满足定理假设之一。(b) 定理 8.7 对 $\mu_0$ 或 $\mu_1$ 绝对连续均成立，故对称。(c) 仍成立，用 $\mu_1\in\Pac$ 的那个方向。}
  % SOLUTION: (a) L=|v|^2满足C^2且Hessian严格正定；C(mu_0,mu_1)<+inf已给；mu_0∈Pac是定理假设之一。定理8.7直接适用，mu_t∈Pac for t∈(0,1)。(b) 定理8.7对"mu_0或mu_1中至少一个绝对连续"均成立，完全对称。(c) 仍成立：用mu_1∈Pac分支，定理8.7的证明对mu_1绝对连续同样有效。

  \TLtakeaway{习题 5.3 训练定理 8.7 假设的精确核查：$L\in C^2$，$\nabla_v^2L>0$，成本有限，端点之一绝对连续——四条缺一不可。}
\end{frame}

% ── §5 小结 ──────────────────────────────────────────────────────────────────
\begin{frame}{§5 小结：无交叉 $\to$ 缩短 $\to$ 内部正则性}
  \textbf{核心思路链：}
  \[
    \text{无交叉}
    \overset{\text{定量化}}{\Rightarrow}
    \text{缩短引理}
    \overset{c\text{-单调}}{\Rightarrow}
    T_{t_0\to t}\text{ Lipschitz}
    \overset{\text{Lusin N}}{\Rightarrow}
    \mu_t\in\Pac
  \]

  \textbf{三个主要结果：}
  \begin{enumerate}\small
    \item \textbf{推论 8.2（缩短引理）：}
      $\sup_t d\le\frac{C_K}{\min(t_0,1-t_0)}\cdot d(t_0)$（$C^2$、严格凸 $L$）。
    \item \textbf{定理 8.5（Lipschitz 传输）：}
      中间时刻 $t_0\in(0,1)$ 的传输映射 $T_{t_0\to t}$ Lipschitz，是唯一 Monge 解。
    \item \textbf{定理 8.7（绝对连续性）：}
      端点之一 $\in\Pac\;\Rightarrow\;\mu_t\in\Pac$ 对所有 $t\in(0,1)$ 成立。
  \end{enumerate}

  \textbf{下节预告。} 第 9–10 章以此内部绝对连续性为出发点，研究 Wasserstein 空间的曲率下界和位移凸性。

  \TLtakeaway{第 8 章精华：最优传输的组合最优性（无交叉）$\to$ 分析正则性（Lipschitz 传输、绝对连续插值）。联结了组合、分析与几何三个视角。}
\end{frame}
